\section{Билет 25. Принцип суперпозиции волн. Сложение когерентных волн. Максимумы и минимумы смещения в результирующей волне. Условия возникновения максимума и минимума смещения.}

\begin{definition}
\textbf{Принцип суперпозиции (наложения) волн} гласит: если в среде одновременно распространяются несколько волн, то результирующее смещение любой точки среды в произвольный момент времени равно векторной (геометрической) сумме смещений, вызываемых каждой волной в отдельности. При этом волны распространяются независимо друг от друга, не искажаясь при встрече.
\end{definition}

\begin{definition}
Наибольший практический интерес представляет сложение волн одинаковой частоты ($\omega_1 = \omega_2 = \omega$), распространяющихся в одном направлении и имеющих постоянную во времени разность фаз. Такие волны называются \textbf{когерентными}. Для плоских монохроматических волн уравнения смещений в произвольной точке имеют вид:
\begin{equation}
    \xi_1 = A_1 \cos(\omega t - k x_1), \quad \xi_2 = A_2 \cos(\omega t - k x_2),
\end{equation}
где $k = 2\pi/\lambda$ --- волновое число, $x_1$ и $x_2$ --- расстояния от источников до точки наблюдения.
\end{definition}

\begin{theorem}[Результирующее колебание]
При наложении когерентных волн в данной точке среды возникает гармоническое колебание той же частоты $\omega$:
\begin{equation}
    \xi = \xi_1 + \xi_2 = A \cos(\omega t - \alpha),
\end{equation}
где амплитуда $A$ результирующего колебания определяется разностью фаз складываемых волн $\Delta\varphi = k(x_2 - x_1) = \frac{2\pi}{\lambda}\Delta x$:
\begin{equation}
    A = \sqrt{A_1^2 + A_2^2 + 2A_1 A_2 \cos(\Delta\varphi)}. \label{eq:res_amp_25}
\end{equation}
Явление устойчивого пространственного распределения амплитуды при интерференции когерентных волн называется \textbf{интерференцией}.
\end{theorem}

\begin{property}[Условия максимума и минимума]
Из формулы \eqref{eq:res_amp_25} следуют условия для амплитуды результирующей волны:
\begin{enumerate}
    \item \textbf{Максимум амплитуды (усиление)} достигается, когда $\cos(\Delta\varphi) = 1$, то есть разность фаз кратна $2\pi$:
    \begin{equation}
        \Delta\varphi = 2n\pi \quad (n=0, 1, 2, \dots) \quad \Rightarrow \quad \Delta x = n\lambda. \label{eq:max_condition}
    \end{equation}
    В этом случае $A_{\max} = A_1 + A_2$. Условие \eqref{eq:max_condition} означает, что максимумы возникают в точках, где разность хода волн равна целому числу длин волн.
    \item \textbf{Минимум амплитуды (ослабление)} достигается, когда $\cos(\Delta\varphi) = -1$, то есть разность фаз равна нечётному числу $\pi$:
    \begin{equation}
        \Delta\varphi = (2n+1)\pi \quad (n=0, 1, 2, \dots) \quad \Rightarrow \quad \Delta x = (2n+1)\frac{\lambda}{2}. \label{eq:min_condition}
    \end{equation}
    В этом случае $A_{\min} = |A_1 - A_2|$. Минимумы возникают в точках, где разность хода равна нечётному числу полуволн.
\end{enumerate}
\end{property}

\begin{figure}[htbp]
    \centering

    % === РИСУНОК 1: Интерференция от двух источников ===
    \begin{tikzpicture}[scale=1.2]
        \node[font=\bfseries] at (0, 5) {Интерференция от двух источников};

        % Источники
        \fill[red] (-1.5, 0) circle (3pt) node[below] {$S_1$};
        \fill[red] (1.5, 0) circle (3pt) node[below] {$S_2$};

        % Волновые фронты от S1 (синие окружности)
        \foreach \r in {1, 2, 3, 4} {
            \draw[blue, thin] (-1.5, 0) circle (\r);
        }

        % Волновые фронты от S2 (красные окружности)
        \foreach \r in {1, 2, 3, 4} {
            \draw[red, thin] (1.7, 0) circle (\r);
        }

        % Линия максимумов (центральная)
        \draw[green!60!black, ultra thick] (0, 0) -- (0, 3.5);
        \node[green!60!black, right] at (0, 3) {максимум};

        % Точка наблюдения P
        \fill[black] (2, 2.5) circle (2pt) node[right] {$P$};
        \draw[blue, dashed] (-1.5, 0) -- (2, 2.5) node[midway, above left] {$r_1$};
        \draw[red, dashed] (1.5, 0) -- (2, 2.5) node[midway, above right] {$r_2$};

        % Разность хода
        \node[align=center, draw, fill=yellow!10] at (-3.5, 2)
            {$\Delta x = r_2 - r_1$\\[4pt]
            Максимум: $\Delta x = n\lambda$\\[2pt]
            Минимум: $\Delta x = (2n+1)\frac{\lambda}{2}$};

        % Оси
        \draw[->] (-4, 0) -- (4, 0) node[right] {$x$};
        \draw[->] (0, -0.5) -- (0, 3.5);
    \end{tikzpicture}

    \vspace{1cm}

    % === РИСУНОК 2: Сложение волн ===
    \begin{tikzpicture}[scale=1.2]
        \node[font=\bfseries] at (4, 2.2) {Сложение когерентных волн};

        % Оси
        \draw[->] (-0.5, 0) -- (7, 0) node[right] {$t$};
        \draw[->] (0, -2.5) -- (0, 2.5) node[above] {$\xi$};

        % Волна 1
        \draw[blue, thick, domain=0:6.5, samples=100]
            plot (\x, {sin(\x*60)});
        \node[blue] at (0.5, 2.2) {$\xi_1$};

        % Волна 2 (в фазе с первой)
        \draw[red, dashed, thick, domain=0:6.5, samples=100]
            plot (\x, {0.8*sin(\x*60)});
        \node[red] at (0.5, 1.7) {$\xi_2$};

        % Результат (усиление)
        \draw[green!60!black, thick, domain=0:6.5, samples=100]
            plot (\x, {1.8*sin(\x*60)});
        \node[green!60!black] at (3.5, 1.7) {$\xi = \xi_1 + \xi_2$ (усиление)};

        % Разность фаз
        \draw[<->, purple] (0, -2.2) -- (1.05, -2.2)
            node[midway, below, font=\small] {$\Delta\varphi = 0$};

        % Подпись
        \node[align=center, font=\small] at (7, -1.5)
            {При $\Delta\varphi = 0$:\\амплитуды\\складываются};
    \end{tikzpicture}

    \caption{Интерференция когерентных волн. Вверху: интерференционная картина от двух источников $S_1$ и $S_2$; точки пересечения волновых фронтов соответствуют максимумам. Внизу: сложение волн при $\Delta\varphi = 0$ (усиление).}
    \label{fig:interference}
\end{figure}

